% GENERATED FILE. Edit canonical lessons, then run npm run generate:books.
% canonical-source-hash: fnv1a64:215bbdbf80f109d8
% canonical-lessons: TE-C529-menamama, TE-C529-menatta, TE-C529-peddananna, TE-C529-peddamma, TE-C529-menakodalu

\chapter{Uncles and Aunts}
\label{ch:te-uncles-and-aunts}
\hlchaptermodality{529}

\begin{chapteropening}
\textbf{By the end of this chapter:} \emph{I can name my uncles, aunts and nieces.}

Complete the last lesson of chapter 529: I can name my uncles, aunts and nieces.
\end{chapteropening}

\section[\texorpdfstring{\te{మేనమామ} (mēnamāma)}{మేనమామ (mēnamāma)}]{\te{మేనమామ} (mēnamāma) — a maternal uncle}

\label{lesson:TE-C529-menamama}



\begin{quote}
Before the new one: say the Telugu for a husband's sister, then the Telugu for an elder brother's wife.
\end{quote}

\subsection*{You'll want to know: \te{మేనమామ}}
\textbf{\te{మేనమామ}} — \emph{mēnamāma} — \textquotedblleft{}a maternal uncle\textquotedblright{}.

\begin{quote}
\te{సెలవుల్లో} \te{మేము} \te{మేనమామ} \te{ఇంటికి} \te{వెళ్ళాము}.
\end{quote}

\begin{quote}
\emph{selavullō mēmu mēnamāma iṇṭiki veḷḷāmu.}
\end{quote}

\begin{quote}
In the holidays we went to our maternal uncle's house.
\end{quote}

\te{మామ} alone can mean father-in-law; \te{మేన}- makes it the mother's brother.

\begin{cousinweb}
\textbf{\te{మేనమామ}} (\emph{mēnamāma}) joins native \emph{māma}, \textquotedblleft{}maternal uncle, father-in-law\textquotedblright{} (DEDR 4813), to \emph{mēna-}, which DEDR places in its appendix of words taken from Indo-Aryan (DEDR App. 53), going back through Prakrit kinship words to Sanskrit \emph{maithuna}.
\end{cousinweb}

\begin{grammarlens}[title={What you've built}]
One more word for describing what you lived and felt.
\end{grammarlens}

\subsection*{Your turn}
\begin{itemize}
\raggedright
  \item \emph{Say it:} \emph{mēnamāma}
  \item \emph{Say it:} \emph{mēnamāma}, once more
  \item \emph{Say it:} \emph{selavullō mēmu mēnamāma iṇṭiki veḷḷāmu.} — In the holidays we went to our maternal uncle's house.
  \item \emph{Recall:} say the Telugu for a husband's sister, then the Telugu for an elder brother's wife, then say \emph{mēnamāma} again
\end{itemize}

\begin{tcolorbox}[breakable,colback=teal!4,colframe=teal!35!black,title={Before you move on}]
What does \te{మేనమామ} mean? (\textquotedblleft{}A maternal uncle\textquotedblright{}.) Say it once more.
\end{tcolorbox}
\section[\texorpdfstring{\te{మేనత్త} (mēnatta)}{మేనత్త (mēnatta)}]{\te{మేనత్త} (mēnatta) — a paternal aunt}

\label{lesson:TE-C529-menatta}



\begin{quote}
Before the new one: say the Telugu for an elder brother's wife, then the Telugu for a maternal uncle.
\end{quote}

\subsection*{You'll want to know: \te{మేనత్త}}
\textbf{\te{మేనత్త}} — \emph{mēnatta} — \textquotedblleft{}a paternal aunt\textquotedblright{}.

\begin{quote}
\te{మా} \te{మేనత్త} \te{పిల్లలతో} \te{నేను} \te{ఆడుకునేవాడిని}.
\end{quote}

\begin{quote}
\emph{mā mēnatta pillalatō nēnu āḍukunēvāḍini.}
\end{quote}

\begin{quote}
I used to play with my father's sister's children.
\end{quote}

The father's sister; \te{అత్త} alone also means mother-in-law.

\begin{cousinweb}
\textbf{\te{మేనత్త}} (\emph{mēnatta}) joins native \emph{atta}, \textquotedblleft{}mother-in-law, father's sister\textquotedblright{} (DEDR 142), to \emph{mēna-}, from DEDR's appendix of Indo-Aryan loans (DEDR App. 53).
\end{cousinweb}

\begin{grammarlens}[title={What you've built}]
One more word for describing what you lived and felt.
\end{grammarlens}

\subsection*{Your turn}
\begin{itemize}
\raggedright
  \item \emph{Say it:} \emph{mēnatta}
  \item \emph{Say it:} \emph{mēnatta}, once more
  \item \emph{Say it:} \emph{mā mēnatta pillalatō nēnu āḍukunēvāḍini.} — I used to play with my father's sister's children.
  \item \emph{Recall:} say the Telugu for an elder brother's wife, then the Telugu for a maternal uncle, then say \emph{mēnatta} again
\end{itemize}

\begin{tcolorbox}[breakable,colback=teal!4,colframe=teal!35!black,title={Before you move on}]
What does \te{మేనత్త} mean? (\textquotedblleft{}A paternal aunt\textquotedblright{}.) Say it once more.
\end{tcolorbox}
\section[\texorpdfstring{\te{పెద్దనాన్న} (peddanānna)}{పెద్దనాన్న (peddanānna)}]{\te{పెద్దనాన్న} (peddanānna) — a father's elder brother}

\label{lesson:TE-C529-peddananna}



\begin{quote}
Before the new one: say the Telugu for a maternal uncle, then the Telugu for a paternal aunt.
\end{quote}

\subsection*{You'll want to know: \te{పెద్దనాన్న}}
\textbf{\te{పెద్దనాన్న}} — \emph{peddanānna} — \textquotedblleft{}a father's elder brother\textquotedblright{}.

\begin{quote}
\te{పెద్దనాన్న} \te{మా} \te{కుటుంబానికి} \te{పెద్ద}.
\end{quote}

\begin{quote}
\emph{peddanānna mā kuṭumbāniki pedda.}
\end{quote}

\begin{quote}
My father's elder brother is the head of our family.
\end{quote}

\te{పెద్ద}, elder, + \te{నాన్న}, father.

\begin{cousinweb}
\textbf{\te{పెద్దనాన్న}} (\emph{peddanānna}) is built from native \emph{pedda}, \textquotedblleft{}great, big, elder\textquotedblright{} (DEDR 4411), and \te{నాన్న}, \textquotedblleft{}father\textquotedblright{}.
\end{cousinweb}

\begin{grammarlens}[title={What you've built}]
One more word for describing what you lived and felt.
\end{grammarlens}

\subsection*{Your turn}
\begin{itemize}
\raggedright
  \item \emph{Say it:} \emph{peddanānna}
  \item \emph{Say it:} \emph{peddanānna}, once more
  \item \emph{Say it:} \emph{peddanānna mā kuṭumbāniki pedda.} — My father's elder brother is the head of our family.
  \item \emph{Recall:} say the Telugu for a maternal uncle, then the Telugu for a paternal aunt, then say \emph{peddanānna} again
\end{itemize}

\begin{tcolorbox}[breakable,colback=teal!4,colframe=teal!35!black,title={Before you move on}]
What does \te{పెద్దనాన్న} mean? (\textquotedblleft{}A father's elder brother\textquotedblright{}.) Say it once more.
\end{tcolorbox}
\section[\texorpdfstring{\te{పెద్దమ్మ} (peddamma)}{పెద్దమ్మ (peddamma)}]{\te{పెద్దమ్మ} (peddamma) — a mother's elder sister}

\label{lesson:TE-C529-peddamma}



\begin{quote}
Before the new one: say the Telugu for a paternal aunt, then the Telugu for a father's elder brother.
\end{quote}

\subsection*{You'll want to know: \te{పెద్దమ్మ}}
\textbf{\te{పెద్దమ్మ}} — \emph{peddamma} — \textquotedblleft{}a mother's elder sister\textquotedblright{}.

\begin{quote}
\te{పెద్దమ్మ} \te{మాకు} \te{అమ్మలాగే}.
\end{quote}

\begin{quote}
\emph{peddamma māku ammalāgē.}
\end{quote}

\begin{quote}
My mother's elder sister is like a mother to us.
\end{quote}

Also the wife of \te{పెద్దనాన్న}.

\begin{cousinweb}
\textbf{\te{పెద్దమ్మ}} (\emph{peddamma}) is built from native \emph{pedda}, \textquotedblleft{}great, big, elder\textquotedblright{} (DEDR 4411), and \te{అమ్మ}, \textquotedblleft{}mother\textquotedblright{}.
\end{cousinweb}

\begin{grammarlens}[title={What you've built}]
One more word for describing what you lived and felt.
\end{grammarlens}

\subsection*{Your turn}
\begin{itemize}
\raggedright
  \item \emph{Say it:} \emph{peddamma}
  \item \emph{Say it:} \emph{peddamma}, once more
  \item \emph{Say it:} \emph{peddamma māku ammalāgē.} — My mother's elder sister is like a mother to us.
  \item \emph{Recall:} say the Telugu for a paternal aunt, then the Telugu for a father's elder brother, then say \emph{peddamma} again
\end{itemize}

\begin{tcolorbox}[breakable,colback=teal!4,colframe=teal!35!black,title={Before you move on}]
What does \te{పెద్దమ్మ} mean? (\textquotedblleft{}A mother's elder sister\textquotedblright{}.) Say it once more.
\end{tcolorbox}
\section[\texorpdfstring{\te{మేనకోడలు} (mēnakōḍalu)}{మేనకోడలు (mēnakōḍalu)}]{\te{మేనకోడలు} (mēnakōḍalu) — a niece; a sister's daughter}

\label{lesson:TE-C529-menakodalu}



\begin{quote}
Before the new one: say the Telugu for a father's elder brother, then the Telugu for a mother's elder sister.
\end{quote}

\subsection*{You'll want to know: \te{మేనకోడలు}}
\textbf{\te{మేనకోడలు}} — \emph{mēnakōḍalu} — \textquotedblleft{}a niece; a sister's daughter\textquotedblright{}.

\begin{quote}
\te{నా} \te{మేనకోడలు} \te{ఇవాళ} \te{బడికి} \te{వెళ్ళింది}.
\end{quote}

\begin{quote}
\emph{nā mēnakōḍalu ivāḷa baḍiki veḷḷindi.}
\end{quote}

\begin{quote}
My niece went to school today.
\end{quote}

A man's sister's daughter or a woman's brother's daughter; \te{కోడలు} alone is a daughter-in-law.

\begin{cousinweb}
\textbf{\te{మేనకోడలు}} (\emph{mēnakōḍalu}) is built from \emph{mēna-}, listed in DEDR's appendix of Indo-Aryan loans (DEDR App. 53) and \te{కోడలు}, \textquotedblleft{}daughter-in-law\textquotedblright{}.
\end{cousinweb}

\begin{grammarlens}[title={What you've built}]
One more word for describing what you lived and felt.
\end{grammarlens}

\subsection*{Your turn}
\begin{itemize}
\raggedright
  \item \emph{Say it:} \emph{mēnakōḍalu}
  \item \emph{Say it:} \emph{mēnakōḍalu}, once more
  \item \emph{Say it:} \emph{nā mēnakōḍalu ivāḷa baḍiki veḷḷindi.} — My niece went to school today.
  \item \emph{Recall:} say the Telugu for a father's elder brother, then the Telugu for a mother's elder sister, then say \emph{mēnakōḍalu} again
\end{itemize}

\begin{tcolorbox}[breakable,colback=teal!4,colframe=teal!35!black,title={Before you move on}]
What does \te{మేనకోడలు} mean? (\textquotedblleft{}A niece; a sister's daughter\textquotedblright{}.) Say it once more.
\end{tcolorbox}
